\chapter{Introdução}

\section{Apresentação}

Esta apostila são notas de aula do curso ministrado por mim no segundo semestre de 2021 para as duas turmas noturnas de Sistemas de Informação ministrado à distância por conta da pandemia de COVID19.
Seu conteúdo está baseado nos livros que servem de bibliografia para o curso:

\nocite{sedgewick01,sedgewick11,cormen12}


\printbibliography[heading=none,keyword={bibliografia}]

%\begin{thebibliography}
%\bibitem{sedgewick11}
%\bibitem{sedgewick01}
%\bibitem{cormen12}
%\end{thebibliography}

  Além dos livros serviu de material para elaboração dessas notas outros textos citados ao longo do curso bem como as gravações dos cursos do professor Robert Sedgeweick para o Coursera e do professor Ronald Rivest para o MIT.

  
  A última versão desta apostila está disponível em um repositório no github (https://github.com/marciomr/apostila-iaa.git).
  Agradeço de antemão eventuais sugestões de correção que forem submetidas pela plataforma.

  Alguns dos direitos sobre o conteúdo desta apostila estão protegidos pela licensa Creative Commons Attribution-NonCommercial-ShareAlike 4.0 International (CC BY-NC-SA 4.0).
Ou seja, você  é livre para distribuir cópias e adaptar este trabalho desde que mantenha a mesma licença, dê o devido crédito ao autor e não faça uso comercial.

\begin{center}
  \includegraphics[width=.3\textwidth]{imagens/cc.png}
\end{center}

\section{Algoritmos}

Um {\em problema computacional} é a especificação de uma relação desejada entre um certo {\em valor de entrada} escolhido em um conjunto de valores válidos e o {\em valor de saída} esperado.

\begin{example}
  {\bf Problema da busca}\\

  {\bf Entrada:} Uma sequência de $n \in \mathbb{N}$ valores $\langle a_1, \dots, a_n \rangle$ em que $a_i \in \mathbb{Z}$ para $1 \leq i \leq n$ e $b \in \mathbb{Z}$.\\

  {\bf Saída:} $i \in \mathbb{N}$ tal que $a_i = b$ se existir ou $\bot$ caso contrário.
\end{example}

\begin{example}
  {\bf Problema da 3-soma}\\

  {\bf Entrada:} Três sequência de $n \in \mathbb{N}$ valores cada $\langle a_1, \dots, a_n \rangle$,  $\langle b_1, \dots, b_n \rangle$ e $\langle c_1, \dots, c_n \rangle$ em que $a_i, b_i, c_i \in \mathbb{Z}$ para $1 \leq i \leq n$.\\

  {\bf Saída:} O número $m \in \mathbb{N}$ de $i, j, k \in \mathbb{N}$ tal que $a_i + b_j + c_j = 0$.
\end{example}

\begin{example}
  {\bf Problema da ordenação}\\

  {\bf Entrada:} Uma sequência de $n \in \mathbb{N}$ valores $\langle a_1, \dots, a_n \rangle$ em que $a_i \in \mathbb{Z}$ para $1 \leq i \leq n$.\\

  {\bf Saída:} Uma permutação da sequência de entrada $\langle a'_1, \dots, a'_n \rangle$ tal que $a_i \leq a_j$ para todo $i \leq j$.
\end{example}

Uma {\em instância do problema} consiste de uma entrada válida para a qual se pretende produzir uma saída esperada.
Por exemplo, a sequência $\langle 3, 42, 17, 2, -1 \rangle$ é uma instância do problema da ordenação cuja saída esperada é $\langle -1, 2, 3, 17, 42 \rangle$.

A disciplina de Introdução à Teoria da Computação (ITC) tem como objeto de estudo problemas computacionais, como eles se classificam entre os que tem solução ou não e entre os que tem solução eficiente ou não.
A solução de um problema computacional é um algoritmo.

Os objetos de estudo desta disciplina são os algoritmos.
Mas afinal, o que são algoritmos?

Um {\em algoritmo} parte de uma entrada escolhida em um conjunto potencialmente infinito de possibilidades ({\em princípio da massividade}) para produzir um valor de saída.
O algoritmo processa a entrada por meio de uma sequência de passos ({\em princípio da discretude}) que produzem valores intermediários.
Cada passo  segue uma instrução simples ({\em princípio da elementaridade}) que só depende dos valores anteriores, não admitindo ambiguidades ({\em princípio da exatidão}) \cite{malcev70}.

Um {\em prorgrama} é a realização de um algoritmo em certa {\em linguagem de programação}.
Assim, um algoritmo é, de um lado, a solução de um problema de computação e, de outro, uma abstração de um conjunto de programas, ele é a idéia por trás do programa.

Um algoritmo é {\em correto} se para toda instância do problema ele produz a saída esperada depois de uma sequência finita de passos.
Nesse caso dizemos que o algoritmo {\em resolve} o problema.

Há uma controversa se devemos ou não considerar uma sequência infinita de instruções como um algoritmo.
Essa questão, complicada, está no coração do nascimento da ciência da computação e será tratada em ITC.
Neste curso focaremos nos algoritmos corretos e, assim, escaparemos da questão.

Para enfatizar o fato de que um algoritmo é uma abstração de um programa, eles serão apresentados neste curso em uma linguagem informal conhecida como {\em pseudo-código}.

\begin{example}
  Considere a seguinte solução para o problema da busca.

\begin{codebox}
\Procname{$\proc{BuscaSimples}(A, b)$}
\li \Comment Recebe $\langle a_1, \dots, a_n \rangle \in \mathbb{Z^n}$ e $b \in \mathbb{Z}$
\li \Comment Devolve $i$ tal que $a_i = b$ se existir e $\bot$ caso contrário
\li \For $i \gets 1$ até $n$
\li \Do \If $a_i = b$
\li     \Then \Return $i$
        \End
    \End
\li \Return $\bot$
\End
\end{codebox}
  
\end{example}

Esta disciplina estuda algoritmos.
Como podemos garantir que certo algoritmo resolve um problema, ou seja, que ele é correto?
Como podemos comparar duas soluções distintas para um mesmo problema?
Ou seja, se conhecemos dois um mais algoritmos corretos para um mesmo problema, como avaliamos qual é melhor?

Avaliaremos os algoritmos corretos a partir da quantidade de recursos que ele consome.
Estudaremos particularmente dois recursos: espaço de memória e, principalmente, o tempo de execução.

Nos capítulos seguintes veremos uma série de algoritmos para resolver alguns problemas centrais da computação como o problema da busca e da ordenação.
Avaliaremos os algoritmos apresentado quanto sua corretude e sua eficiência em consumo de tempo e espaço de memória.

No Capítulo X apresentaremos o estudo dos algoritmos a partir do método empírico.
Relembraremos o método e veremos um exemplo comparando o tempo de execução de duas soluções para o problema da busca em sequências ordenadas.
Então exploraremos técnicas para arrsicar modelos matemáticos adequados para avaliar o consumo de tempo dos algoritmos.
E finalmente veremos ferramentas matemáticas uteis para comparar funções quanto ao seu crescimento, a chamada notação assintótica.
No Capitulo X estudaremos algoritmos de ordenação como estudo de caso da teoria apresentada anteriormente.
Veremos uma série de algoritmos que resolvem o mesmo problema e utilizaremos as técnicas apresentadas para construir e testar modelos do consumo de tempo deles.
Estudaremos também um limite teórico da eficiência do problema da ordenação e veremos dois algoritmos que superam esse limite utilizando mais informações do que as assumidas no enunciado do teorema.
Concluiremos a apostila no Capítulo X apresentando algoritmos e técnicas um pouco mais avançãdos como programação dinâmica e análise amortizada.

